% Source: https://artofproblemsolving.com/community/c3h2564110_mathematics
% Stable identity: aops-topic-2564110-q01
\textbf{(a)} $n=\dfrac27(p+n)$, i.e. $7n=2p+2n$, i.e. $5n=2p$.

\textbf{(b)} After the removal there are $p-84$ pennies and $p-84+n$ coins, so $p-84=\dfrac13(p-84+n)$, i.e. $3p-252=p-84+n$, i.e. $n=2p-168$.

\textbf{(c)} From (a), $2p=5n$, so (b) gives $n=5n-168$, i.e. $n=42$ (and $p=105$). Check: $42$ of the $147$ coins are nickels, which is $\dfrac27$; after the removal there are $21$ pennies among $63$ coins, which is $\dfrac13$. There are $\boxed{42}$ nickels.
